% GENERATED FILE. Edit canonical lessons, then run npm run generate:books.
% canonical-source-hash: fnv1a64:5d8abb50b2265c98
% canonical-lessons: GE-C414-anschluss, GE-C414-nahverkehr, GE-C414-fernverkehr, GE-C414-pendeln, GE-C414-pendler

\chapter{Jeden Tag mit dem Zug}
\label{ch:ge-jeden-tag-mit-dem-zug}
\hlchaptermodality{414}

\begin{chapteropening}
\textbf{By the end of this chapter:} \emph{I can talk about travelling to work: the connection, local transport, long-distance transport, commuting and the commuter.}

Complete the last lesson of chapter 414: I can talk about travelling to work: the connection, local transport, long-distance transport, commuting and the commuter.
\end{chapteropening}

\section[der Anschluss]{der Anschluss — the connection (to another train)}

\label{lesson:GE-C414-anschluss}



\begin{quote}
Before the new one: say the German for to book, then the German for to travel on.
\end{quote}

\subsection*{You'll want to know: der Anschluss}
\textbf{der Anschluss} — \textquotedblleft{}the connection (to another train)\textquotedblright{}.

\begin{quote}
Wegen der Verspätung verpassen wir den Anschluss.
\end{quote}

\begin{quote}
\emph{Because of the delay we miss our connection.}
\end{quote}

Plural: \emph{die Anschlüsse}. Also a connection for water, gas or the internet.

\begin{cousinweb}
\emph{Anschluss} is built on \emph{anschließen} 'to connect, to join on' (EtymDB), from \emph{schließen} 'to close, to lock', inherited from Old High German \emph{sliozan} (EtymDB).
\end{cousinweb}

\begin{grammarlens}[title={What you've built}]
One more word for telling, explaining and describing at B1.
\end{grammarlens}

\subsection*{Your turn}
\begin{itemize}
\raggedright
  \item \emph{Say it:} \emph{der Anschluss}
  \item \emph{Say it:} \emph{der Anschluss}, once more
  \item \emph{Say it:} a sentence of your own with \emph{der Anschluss}
  \item \emph{Recall:} say the German for to book, then the German for to travel on, then say \emph{der Anschluss} again
\end{itemize}

\begin{tcolorbox}[breakable,colback=teal!4,colframe=teal!35!black,title={Before you move on}]
What does der Anschluss mean? (\textquotedblleft{}The connection (to another train)\textquotedblright{}.) Say it once more.
\end{tcolorbox}
\section[der Nahverkehr]{der Nahverkehr — the local transport}

\label{lesson:GE-C414-nahverkehr}



\begin{quote}
Before the new one: say the German for to travel on, then the German for the connection.
\end{quote}

\subsection*{You'll want to know: der Nahverkehr}
\textbf{der Nahverkehr} — \textquotedblleft{}the local transport\textquotedblright{}.

\begin{quote}
Mit dem Nahverkehr kommt man gut ins Zentrum.
\end{quote}

\begin{quote}
\emph{With local transport you can get to the centre easily.}
\end{quote}

No plural. \emph{Der öffentliche Nahverkehr} — public transport in a town and its region.

\begin{cousinweb}
\emph{Nahverkehr} is built from \emph{nah} 'near', inherited from Old High German \emph{nāh} and a cousin of English \emph{nigh} (EtymDB), and \emph{Verkehr} 'traffic', from \emph{kehren} 'to turn'.
\end{cousinweb}

\begin{grammarlens}[title={What you've built}]
One more word for telling, explaining and describing at B1.
\end{grammarlens}

\subsection*{Your turn}
\begin{itemize}
\raggedright
  \item \emph{Say it:} \emph{der Nahverkehr}
  \item \emph{Say it:} \emph{der Nahverkehr}, once more
  \item \emph{Say it:} a sentence of your own with \emph{der Nahverkehr}
  \item \emph{Recall:} say the German for to travel on, then the German for the connection, then say \emph{der Nahverkehr} again
\end{itemize}

\begin{tcolorbox}[breakable,colback=teal!4,colframe=teal!35!black,title={Before you move on}]
What does der Nahverkehr mean? (\textquotedblleft{}The local transport\textquotedblright{}.) Say it once more.
\end{tcolorbox}
\section[der Fernverkehr]{der Fernverkehr — the long-distance transport}

\label{lesson:GE-C414-fernverkehr}



\begin{quote}
Before the new one: say the German for the connection, then the German for the local transport.
\end{quote}

\subsection*{You'll want to know: der Fernverkehr}
\textbf{der Fernverkehr} — \textquotedblleft{}the long-distance transport\textquotedblright{}.

\begin{quote}
Im Fernverkehr sind die Züge schneller und teurer.
\end{quote}

\begin{quote}
\emph{On long-distance routes the trains are faster and dearer.}
\end{quote}

No plural. \emph{Der Fernzug} — the long-distance train.

\begin{cousinweb}
\emph{Fernverkehr} is built from \emph{fern} 'far', inherited from Old High German \emph{ferrana} 'from afar' (EtymDB), and \emph{Verkehr} 'traffic', from \emph{kehren} 'to turn'.
\end{cousinweb}

\begin{grammarlens}[title={What you've built}]
One more word for telling, explaining and describing at B1.
\end{grammarlens}

\subsection*{Your turn}
\begin{itemize}
\raggedright
  \item \emph{Say it:} \emph{der Fernverkehr}
  \item \emph{Say it:} \emph{der Fernverkehr}, once more
  \item \emph{Say it:} a sentence of your own with \emph{der Fernverkehr}
  \item \emph{Recall:} say the German for the connection, then the German for the local transport, then say \emph{der Fernverkehr} again
\end{itemize}

\begin{tcolorbox}[breakable,colback=teal!4,colframe=teal!35!black,title={Before you move on}]
What does der Fernverkehr mean? (\textquotedblleft{}The long-distance transport\textquotedblright{}.) Say it once more.
\end{tcolorbox}
\section[pendeln]{pendeln — to commute; to swing}

\label{lesson:GE-C414-pendeln}



\begin{quote}
Before the new one: say the German for the local transport, then the German for the long-distance transport.
\end{quote}

\subsection*{You'll want to know: pendeln}
\textbf{pendeln} — \textquotedblleft{}to commute; to swing\textquotedblright{}.

\begin{quote}
Sie pendelt jeden Tag zwischen zwei Städten.
\end{quote}

\begin{quote}
\emph{She commutes between two towns every day.}
\end{quote}

Perfekt with \emph{sein} for travel: \emph{er ist jahrelang gependelt}.

\begin{cousinweb}
\emph{pendeln} is built on \emph{Pendel} 'pendulum' (EtymDB), a learned borrowing from Latin \emph{pendulus} 'hanging', from \emph{pendēre} 'to hang' (Lewis \& Short): a commuter swings back and forth like a pendulum.
\end{cousinweb}

\begin{grammarlens}[title={What you've built}]
One more word for telling, explaining and describing at B1.
\end{grammarlens}

\subsection*{Your turn}
\begin{itemize}
\raggedright
  \item \emph{Say it:} \emph{pendeln}
  \item \emph{Say it:} \emph{pendeln}, once more
  \item \emph{Say it:} a sentence of your own with \emph{pendeln}
  \item \emph{Recall:} say the German for the local transport, then the German for the long-distance transport, then say \emph{pendeln} again
\end{itemize}

\begin{tcolorbox}[breakable,colback=teal!4,colframe=teal!35!black,title={Before you move on}]
What does pendeln mean? (\textquotedblleft{}To commute; to swing\textquotedblright{}.) Say it once more.
\end{tcolorbox}
\section[der Pendler]{der Pendler — the commuter}

\label{lesson:GE-C414-pendler}



\begin{quote}
Before the new one: say the German for the long-distance transport, then the German for to commute.
\end{quote}

\subsection*{You'll want to know: der Pendler}
\textbf{der Pendler} — \textquotedblleft{}the commuter\textquotedblright{}.

\begin{quote}
Am Morgen ist der Zug voll mit Pendlern.
\end{quote}

\begin{quote}
\emph{In the morning the train is full of commuters.}
\end{quote}

Same form in the plural. A woman: \emph{die Pendlerin}.

\begin{cousinweb}
\emph{Pendler} is built on \emph{pendeln} 'to commute' (EtymDB), from \emph{Pendel} 'pendulum', which goes back to Latin \emph{pendulus} 'hanging' (Lewis \& Short).
\end{cousinweb}

\begin{grammarlens}[title={What you've built}]
One more word for telling, explaining and describing at B1.
\end{grammarlens}

\subsection*{Your turn}
\begin{itemize}
\raggedright
  \item \emph{Say it:} \emph{der Pendler}
  \item \emph{Say it:} \emph{der Pendler}, once more
  \item \emph{Say it:} a sentence of your own with \emph{der Pendler}
  \item \emph{Recall:} say the German for the long-distance transport, then the German for to commute, then say \emph{der Pendler} again
\end{itemize}

\begin{tcolorbox}[breakable,colback=teal!4,colframe=teal!35!black,title={Before you move on}]
What does der Pendler mean? (\textquotedblleft{}The commuter\textquotedblright{}.) Say it once more.
\end{tcolorbox}
